% =============================================================================
%  SPECTRA — Section I: Introduction (body) + Section VIII: Conclusion
%  Target: IEEE Transactions on Information Forensics & Security
%  Author: Sagar Korde
%  Note: This file contains the Introduction body (after opening paragraphs)
%        and the Conclusion section. Merge with master .tex accordingly.
% =============================================================================

% -----------------------------------------------------------------------------
%  SECTION I — INTRODUCTION (body, to be placed after opening paragraphs)
% -----------------------------------------------------------------------------

\section{Introduction}
\label{sec:introduction}

Bitcoin's pseudonymous ledger now records more than \num{5.8} million
transactions daily, creating an enormous evidentiary surface for forensic
investigators, compliance teams, and anti-money-laundering (AML) practitioners.
Extracting actionable intelligence from this stream requires solving three
structurally distinct problems simultaneously.
First, \emph{transaction-type identification}: determining whether a given
transaction represents an ordinary peer-to-peer transfer, a coin-join mixing
event, a service payment, or another category, is essential for financial
forensics yet remains challenging because transaction types share overlapping
fee, size, and graph-degree distributions.
Second, \emph{real-time anomaly and trust scoring}: AML compliance mandates
that suspicious clusters of addresses be flagged with quantified confidence,
not merely a binary label, and that these scores update continuously as new
blocks arrive.
Third, \emph{temporal concept drift}: the Bitcoin ecosystem evolves---protocol
upgrades (SegWit, Taproot), shifting exchange policies, and the entry of new
institutional actors cause the statistical distribution of transactions to
shift over time in ways that silently degrade frozen classifiers without
triggering any observable runtime error.
Existing methods address at most one of these challenges.
Rule-based heuristics~\cite{Harrigan2016} provide labels but no uncertainty
estimates and become stale as the protocol evolves.
Supervised classifiers trained on static, rule-annotated datasets achieve
high accuracy on held-out splits from the same distribution but generalize
poorly when the annotation rules themselves constitute features exploited by
the model.
Graph neural networks proposed for blockchain analysis
\cite{Pareja2020,Weber2019} improve structural reasoning but offer no
probabilistic guarantees, no drift monitoring, and no per-address trust
certificates.
No single published framework simultaneously addresses all three challenges
within a reproducible, end-to-end pipeline.

\subsection*{Gap Analysis}

The dominant family of tabular classifiers---gradient boosted
trees~\cite{Chen2016} being the canonical example---achieves near-perfect
accuracy on rule-annotated blockchain benchmarks.
This performance, however, reflects a form of \emph{label--feature
entanglement}: because transaction-type labels are themselves derived from
deterministic heuristics applied to the same raw features that the classifier
observes (e.g., the number of inputs and outputs, the fee-per-byte ratio),
the model learns to invert the labelling function rather than to reason about
latent economic behavior.
When temporal identifiers such as \texttt{block\_height}, \texttt{block\_time},
and \texttt{year} are included---as they are in standard benchmark releases
---a decision tree can partition the data by era and recover the exact rule
boundaries without any generalization whatsoever.
XGBoost and BERT4ETH~\cite{Hu2023} achieve accuracy of \num{1.0} under
this oracle setting; removing these temporally entangled features degrades
their accuracy by more than \num{30} percentage points, exposing the
fragility of this evaluation paradigm.

Graph-based approaches address topology but introduce their own limitations.
EvolveGCN~\cite{Pareja2020} adapts GCN weights across temporal snapshots
but treats each snapshot independently, cannot propagate trust estimates,
and provides no mechanism for detecting the distributional shift that triggers
its own accuracy degradation.
AA-GNN~\cite{Zhou2020} augments attention over neighborhood aggregation but
similarly lacks probabilistic uncertainty quantification.
Clustering-based methods~\cite{McInnes2017} identify behavioral cohorts in
feature space without providing classification or drift certificates.
Bayesian methods applicable in principle to address scoring
\cite{DeFord2019} have not been integrated with graph-structured input
representations.
To the best of our knowledge, no prior work simultaneously provides
(i)~density-based behavioral clustering in a learned latent space,
(ii)~drift monitoring via optimal-transport distances,
(iii)~per-address Bayesian trust certificates, and
(iv)~inductive node classification over the transaction graph.

\subsection*{SPECTRA: Overview}

We present \textsc{SPECTRA}
(\textbf{S}pectral-\textbf{P}robabilistic \textbf{E}nsemble \textbf{C}lustering
for \textbf{T}ransaction \textbf{R}ecognition and \textbf{A}uthentication),
a seven-module framework designed to address all three challenges within a
unified, interpretable pipeline.

\textbf{Module S} constructs a directed weighted address graph
$\mathcal{G} = (\mathcal{V}, \mathcal{E}, \mathbf{W})$ over \num{30000}
high-degree hub nodes and extracts $k=10$ spectral node embeddings from the
normalized graph Laplacian via randomized truncated SVD~\cite{Halko2011}.
Randomized SVD is critical at Bitcoin scale: the ARPACK eigensolver, which
underlies SciPy's \texttt{eigsh}, suffers convergence failures and threading
crashes on sparse graphs exceeding $10^6$ nodes under Python~3.12's GIL
implementation; randomized methods achieve $\epsilon$-approximate
eigenvectors in $\mathcal{O}(nk\log k)$ time regardless of graph size.
The resulting embeddings capture global community structure in a form stable
across block-height partitions.

\textbf{Modules E and C} form the probabilistic core of the framework.
Module~E applies mutual-information-based feature selection followed by
non-negative matrix factorization to identify the top-20 latent transaction
attributes.
Module~C encodes each transaction into a Gaussian latent distribution
$q_\phi(\mathbf{z}|\mathbf{x}) = \mathcal{N}(\boldsymbol{\mu}, \boldsymbol{\sigma}^2\mathbf{I})$
via a $\beta$-VAE~\cite{Kingma2014} with latent dimension \num{32},
enabling principled anomaly scoring through the Evidence Lower Bound (ELBO)
reconstruction error.
HDBSCAN~\cite{McInnes2017} applied to this latent space discovers \num{135}
behavioral clusters without requiring a pre-specified cluster count $k$,
grouping addresses by their transaction patterns rather than by arbitrary
geometric boundaries.

\textbf{Module C} additionally provides temporal drift monitoring: Wasserstein-1
distances between successive cluster-label distributions are computed over
rolling time windows of configurable width.
Across \num{29} of \num{29} evaluated windows in our experimental protocol,
the Wasserstein distance exceeds the alert threshold $\delta = 0.1$,
confirming that Bitcoin transaction behavior exhibits persistent, statistically
significant concept drift---a finding with direct implications for the
deployment lifetime of any frozen classifier.

\textbf{Modules P and T} produce address-level intelligence.
Module~P maintains per-address Bayesian trust scores modeled as a Beta
conjugate posterior, updated online from reconstruction error likelihoods,
with a dataset-wide mean trust of \num{0.634}.
Module~T constructs Markov chain behavioral profiles over the \num{136}
observed cluster states, scoring anomaly via the negative log-probability
of an address's observed transition sequence; the mean anomaly score of
\num{0.875} reflects the high entropy of transaction behavior in the
CryptoGraphNet corpus.

\textbf{Module R} fuses all preceding module outputs into a node-feature
matrix and trains a two-layer GraphSAGE network~\cite{Hamilton2017} for
inductive address classification.
GraphSAGE's sampling-based neighborhood aggregation generalizes to addresses
unseen during training, a property essential for live blockchain monitoring
where new addresses are created at each block.

\textbf{Module A} integrates anomaly signals from the VAE (ELBO), Markov
chain (transition probability), and Bayesian trust score into a unified
per-address anomaly index, providing a single, interpretable output for
AML alert triage.

\subsection*{Evaluation Preview}

We evaluate SPECTRA on the CryptoGraphNet
benchmark~\cite{Farrugia2020} comprising \num{300000} transactions
labeled across six transaction types.
To address the label--feature entanglement problem described above, we
introduce a \emph{two-track evaluation protocol}: a \emph{tabular track}
that evaluates classifiers with access to all raw features (including temporal
identifiers), and a \emph{graph-topology track} that restricts all models
to graph-structural and module-derived features only.
Under this protocol, XGBoost and BERT4ETH achieve accuracy of \num{1.0} on
the tabular track (oracle performance attributable to entanglement) but
degrade substantially on the graph-topology track, while SPECTRA achieves
accuracy of \num{0.916} and Matthews Correlation Coefficient (MCC) of
\num{0.840}---outperforming EvolveGCN by \num{33.9} percentage points and
AA-GNN by \num{41.7} percentage points on the same track.
Across six ablation variants, the full seven-module system consistently
dominates all partial ablations, confirming that each module contributes
independently to classification performance.

\subsection*{Contributions}

The primary contributions of this paper are as follows:

\begin{enumerate}[leftmargin=*, label=\textbf{C\arabic*}]

  \item \textbf{Spectral graph representation with randomized SVD
    (Module~S).}
    We construct a directed weighted Bitcoin address graph and extract
    $k$-dimensional spectral embeddings from the normalized graph Laplacian
    via randomized truncated SVD~\cite{Halko2011}, resolving convergence and
    threading failures encountered by ARPACK on large sparse graphs and
    providing provably convergent spectral features for nodes of any
    connectivity degree.

  \item \textbf{Probabilistic VAE latent space with principled anomaly
    scoring (Modules~E/C).}
    A $\beta$-VAE~\cite{Kingma2014} encodes transactions into a
    \num{32}-dimensional Gaussian latent space, enabling ELBO-based anomaly
    scoring and density clustering (HDBSCAN~\cite{McInnes2017}) that
    discovers \num{135} behavioral clusters without a pre-specified $k$.

  \item \textbf{Wasserstein drift monitoring over cluster distributions
    (Module~C).}
    We compute Wasserstein-1 distances between successive cluster-label
    distributions over rolling time windows and provide quantitative,
    threshold-based drift alerts, detecting persistent concept drift in
    \num{29} of \num{29} evaluation windows.

  \item \textbf{Bayesian trust certificates and Markov behavioral profiling
    (Modules~P/T).}
    Per-address trust is maintained as a Beta conjugate posterior updated
    online from reconstruction error likelihoods.
    A Markov chain over \num{136} cluster states captures behavioral
    trajectories, with anomaly scores derived from low-probability
    state transitions.

  \item \textbf{Inductive GraphSAGE classifier (Module~R).}
    A two-layer GraphSAGE network~\cite{Hamilton2017} trained on
    module-fused node features performs address classification inductively,
    generalizing to addresses unseen during training---a critical property
    for live blockchain monitoring.

  \item \textbf{Temporally-fair two-track evaluation protocol.}
    We identify and formally characterize label--feature entanglement in
    rule-annotated blockchain datasets, propose a two-track evaluation
    protocol separating tabular and graph-topology methods, and release
    temporally-fair baseline results with \texttt{block\_height},
    \texttt{block\_time}, and \texttt{year} removed to prevent temporal
    identity exploitation.

  \item \textbf{Comprehensive, reproducible benchmark.}
    We evaluate ten baselines (K-Means, DBSCAN, Isolation Forest,
    AE+KMeans, HDBSCAN, XGBoost, EvolveGCN, BERT4ETH-BTC, Graph
    Autoencoder, AA-GNN) and six ablation variants under identical
    chronological train/validation/test splits, with all code, configurations,
    and checkpoints publicly released.

\end{enumerate}

The remainder of this paper is organized as follows.
Section~\ref{sec:related} surveys related work.
Section~\ref{sec:problem} formalizes the problem.
Section~\ref{sec:architecture} describes the SPECTRA architecture in detail.
Section~\ref{sec:label_provenance} analyzes label--feature entanglement.
Section~\ref{sec:experiments} presents the experimental setup.
Section~\ref{sec:results} reports results and ablations.
Section~\ref{sec:conclusion} concludes with limitations and future directions.


% -----------------------------------------------------------------------------
%  SECTION VIII — CONCLUSION
% -----------------------------------------------------------------------------

\section{Conclusion}
\label{sec:conclusion}

We have presented \textsc{SPECTRA}, a seven-module framework for Bitcoin
transaction analysis that, to the best of our knowledge, is the first system
to simultaneously address transaction-type classification, per-address
Bayesian trust scoring, temporal concept-drift monitoring, and interpretable
behavioral clustering within a single end-to-end pipeline.
On the CryptoGraphNet benchmark, SPECTRA-full achieves an accuracy of
\num{0.916} and a Matthews Correlation Coefficient of \num{0.840} on the
graph-topology evaluation track, outperforming EvolveGCN~\cite{Pareja2020}
by \num{33.9} percentage points and AA-GNN~\cite{Zhou2020} by
\num{41.7} percentage points.
Beyond classification, the framework discovers \num{135} behavioral
clusters without a pre-specified $k$, detects statistically significant
concept drift in \num{29} of \num{29} evaluated time windows, and assigns
per-address Bayesian trust scores with a dataset-wide mean of \num{0.634}
---outputs that no single prior baseline provides in isolation.

A methodological contribution of equal practical significance is our
two-track evaluation protocol for label-entangled datasets.
By separating the tabular track (where rule-derived features allow oracle-level
accuracy) from the graph-topology track (where only structural and
module-derived features are permitted), we expose the fragility of
previously reported benchmark results and provide a reproducible framework
for fair comparison in future work.
We recommend that all future evaluations on rule-annotated blockchain
benchmarks adopt this distinction explicitly.

\subsection*{Limitations}

Three limitations warrant disclosure.
First, \emph{hub-node coverage}: the spectral graph module operates on the
\num{30000} highest-degree addresses; \num{78.5}\% of test transactions
involve addresses absent from this hub set and therefore receive a uniform
prior rather than a learned spectral embedding.
Bitcoin's one-time-address model---where addresses are routinely discarded
after a single use---makes this an inherent challenge for any graph-based
approach that requires sufficient structural context per node.
Second, \emph{HDBSCAN memory complexity}: HDBSCAN's mutual-reachability
construction is $\mathcal{O}(n^2)$ in memory for exact computation, which
limits scalability to the $10^6$-transaction regime without approximate
nearest-neighbor pre-filtering.
Third, \emph{Markov chain order assumption}: Module~T assumes the first-order
Markov property---that the next cluster state depends only on the current
state---which may underfit address sequences with longer-range behavioral
dependencies such as multi-round coin-join protocols.

\subsection*{Future Work}

Four directions present themselves as natural extensions.
\textbf{Cross-chain generalization}: applying SPECTRA to Ethereum and
Monero requires handling account-model semantics and ring-signature
unlinkability, respectively, and provides a stringent test of the
framework's architectural assumptions.
\textbf{Federated forensics}: a federated learning variant of the
GraphSAGE and VAE modules would allow multiple exchanges or law-enforcement
agencies to collaboratively train shared representations without disclosing
raw transaction data---a privacy-preserving configuration aligned with
GDPR and FATF guidelines for multi-party AML compliance.
\textbf{Attention-based behavioral profiling}: replacing the first-order
Markov chain in Module~T with a transformer-based sequence model or a
higher-order Markov process would capture multi-hop behavioral dependencies
and likely improve anomaly detection precision for sophisticated mixing
schemes.
\textbf{Online GNN fine-tuning for address drift}: as new address classes
emerge with each protocol upgrade, a continual-learning extension of
Module~R that fine-tunes GraphSAGE weights online---without catastrophic
forgetting of previously learned classes---would extend the operational
lifetime of a deployed SPECTRA instance without requiring full retraining.

\textsc{SPECTRA}'s modular architecture is designed to accommodate these
extensions without redesign of the core pipeline, and all components are
released to facilitate community development.
